The computational model is C++17. It uses the specified 250 m floodplain
grid and rectangular channel sections at approximately 250 m spacing.
The channel mask consists of grid centers within 250 m of the centerline.
The nearest segment gives linear interpolation weights, which are also used
for incoming exchange. An outgoing mass is removed in proportion to each
section's contribution to the interpolated donor inventory; this prevents
withdrawal from an empty section. The channel's exact rectangular control
areas replace masked floodplain storage. The raster bed inventory and the
rectangular water storage consequently have slightly different footprints.

We inherit channel and floodplain Chezy coefficients of 55 and 10
m$^{1/2}$ s$^{-1}$ from Nicholas et al. (2013) \cite{nicholas2013}. Let $C_h$ denote this
coefficient, distinct from sediment coefficient $C$. The resistance and
eddy viscosity are
\begin{equation}
 \mathbf f_b=\frac{g}{C_h^2}\mathbf u|\mathbf u|,\qquad
 u_*=\frac{\sqrt g}{C_h}|\mathbf u|,\qquad \nu_t=0.07u_*h.
\end{equation}
The symmetric velocity-gradient tensor gives horizontal turbulent stress;
we set sediment diffusivity equal to $\nu_t$, an assumed Schmidt number of
one. A local friction update follows each transport substep.

We assume that $S$ is in mg L$^{-1}$, consistent with concentration units
reported in the 2013 methods paper, so internal inlet concentration is
$10^{-3}CQ^{3/2}$ kg m$^{-3}$. At the peak discharge, $C=0.003$ gives
5.511 kg m$^{-3}$. The 2018 paper does not establish this unit convention.
Porosity 0.4 and grain density 2650 kg m$^{-3}$ are assumed. Sand settles
at 0.023244 m s$^{-1}$, calculated from the Ferguson--Church relation \cite{ferguson2004} for
assumed 0.2 mm particles. Silt uses 0.0002 m s$^{-1}$ from the earlier
methods paper. Clay uses Stokes settling for assumed unflocculated 2 $\mu$m
spheres, giving $3.597\times10^{-6}$ m s$^{-1}$.
For silt and clay, the deposition sink is
\begin{equation}
 d_k=w_{s,k}c_k\max(0,1-\tau_b/\tau_c),\qquad
 \tau_b=\rho_w g|\mathbf u|^2/C_h^2,
\end{equation}
with assumed $\tau_c=0.1$ Pa. Sand deposition is independent of this shear
reduction. At fixed depth and stress, deposition integrates exponentially,
and the removed suspended mass increments bed elevation exactly once.
Vertical floodplain entrainment is omitted because the earlier
sand-capacity-based closure does not close the three-fraction,
prescribed-channel-bed problem. This omission is a physical uncertainty.

The floodplain is initially dry. The channel starts at the minimum-discharge
normal depth, 4.356 m, with zero suspended concentration. The unreported
hydrograph shape is taken as
\begin{equation}
 Q(t)=2500+12500\sin^2(\pi t/T_f)\quad\mathrm{m^3\,s^{-1}}.
\end{equation}
Water and sediment fields persist between floods. At the inlet we prescribe
$Q$ and preserve the outgoing shallow-water characteristic
$u-2\sqrt{gh}$. At the subcritical outlet we preserve $u+2\sqrt{gh}$ and
impose the Chezy normal-flow rating for the current channel slope.
These boundary conditions are assumptions. Unsupported incoming
supercritical states stop the calculation with an explicit error.
Floodplain lateral and upstream boundaries reflect flow; the downstream
boundary is transmissive with continued valley slope.

Floodplain transport rates are evaluated once per macro step, with CFL
coefficient 0.4, a 120 s ceiling, and positive-inventory restrictions.
The channel takes additional wave-limited substeps with fixed lateral
exchange rates during that interval. The scheme is first order in time.
Both representations receive identical integrated water and sediment
exchange; channel momentum exchange includes the advective streamwise
component. Side pressure is normal to the physical streamline and is not
added as a streamwise force from raster stair steps. Exactly dry regions
are skipped computationally without changing active-state arithmetic.

Migration uses the linear-curvature Howard--Knutson/Ikeda parameterization
written in later methods descriptions \cite{schwendel2015,limaye2014}. For curvature $\kappa$, define
\begin{equation}
 R_0=k_l B\kappa,\qquad
 R_1=E\left[-R_0+2.5\frac{\int_0^s R_0(\xi)
 e^{-a(s-\xi)}\,d\xi}{\int_0^s e^{-a(s-\xi)}\,d\xi}\right],
 \quad a=\frac{2g}{C_h^2 H}.
\end{equation}
We use assumed $k_l=400$ m per flood, bankfull depth $H$, fixed endpoints,
and 5 m initial random perturbations with seed 2018. Here $k_l$ is the
coefficient per unit dimensionless curvature, not the mean observed
migration speed. The upstream integral is exact for piecewise-linear
curvature. Adaptive steps and geometric resampling control numerical
crowding. Nonadjacent limbs closer than one channel width form a neck
cutoff. The precise 2018 migration parameterization is unverified.

Migration precedes prescribed aggradation. Newly cut cells are excavated
to the pre-aggradation channel bed; abandoned cells retain their former
bed. Eroded dry mass enters a reservoir and is returned at the next inlet
with the selected grain fractions and a $Q^{3/2}$ concentration shape.
Release is limited by available reservoir inventory, with quadrature
remainders carried forward. To transfer abandoned cells, we normalize the
old interpolation weights over each section's entire raster support. If
$V_jq_j$ is its inventory and $S_j=\sum_r a_rw_{rj}$ its sampled area,
abandoned cell $r$ receives $\sum_j(a_rw_{rj}/S_j)V_jq_j$. The retained
support fraction stays in the channel; unsupported sections retain all their
inventory. Conservative overlap in normalized arc length remaps the remaining
channel state. Newly captured floodplain inventories enter neighboring new
sections using the new interpolation weights. Water and every sediment class
share these nonnegative weights, preserving mass and constant concentrations.
Local depth need not be preserved when raster and rectangular areas differ.
The earlier subtraction of full ghost-cell inventories failed a 250 m cutoff
test by overdrawing a localized sediment class by 35.48\%; that attempt was
stopped and its records retained.
Abandoned momentum rotates into the old segment tangent; incoming floodplain
momentum projects onto the new segment tangent. Geometry change therefore
does not conserve Cartesian momentum. This remapping closure and fixed bank
grain fractions are explicit assumptions. Prescribed bed
changes have a separate signed sediment source in the mass ledger.
A nonpositive bankfull depth is treated as a failure of the prescribed
channel geometry, not replaced by an artificial minimum.

For $\alpha$, the floodplain depth average includes dry cells. A bank face
is inundated when floodplain water and channel overtopping both exceed
1 mm; this threshold affects diagnostics only. No wet interface gives an
undefined value. The immediate active-channel interface is consistent
with visual inspection of published Figure 1E, although the raster
implementation and wetness criterion are assumptions. Ridge relief for
$\beta$ uses the maximum transverse value of longitudinal-mean relief over
the central 80\% of the reach, relative to the outermost 10\% lateral bands.
Channel depth is the bank-face mean elevation minus adjacent channel bed.
We average the resulting $\beta_n$ over floods 561--580. These averaging
choices are not specified in DR1.
